\chapter{Lista 14.}

\begin{context}
	Gdy $R$ jest dziedziną, przez $K$ oznaczamy ciało ułamków $R$, tzn $K = R_0.$
\end{context}
\begin{problem}[1]
	Niech $S$ będzie zbiorem multiplikatywnym $R$. Definiujemy relację $\sim  $ na $R\times S$ przez 
	\[
		\left( r,s \right) \sim \left( r',s' \right) \iff \left(\exists s'' \in S \right) \left(s''\left( rs' -r's \right) = 0\right)   
	.\] 
	Udowodnić, że jest to relacja równoważności. Na $\mfaktor{\left( R\times S \right) }{\sim  } $ definiujemy $\plus$ i  $\cdot $ jak w przypadku lokalizacji dla dziedzin. Udowodnić, że działania te są dobrze zdefiniowane. Pokazać, że spełniają one aksjomaty pierścienia.
\end{problem} 
\begin{solution}
	TODO
\end{solution} 
\begin{problem}[2]
	Załóżmy, że $R$ jest UFD. Niech $\mathcal{P}$ będzie zbiorem reprezentantów klas abstrakcji relacji stowarzyszenia na zbiorze wszystkich elementów nierozkładalnych $R$. Uzasadnić, że każdy element $K\setminus \{0\} $ zapisuje się jednoznacznie (z dokładnością do kolejności czynników) jako produkt $a = \varepsilon \prod_{p \in \mathcal{P}} p^{v_{p}}$, gdzie $\varepsilon \in R^{\ast} $ i $v_{p}\in \mathbb{Z}$.

	Zauważmy, że dzięki temu otrzymaliśmy na wykładzie funkcję $v_{p}: K \to \mathbb{Z}$ (waluacja $p$--adyczna), która posłużyła do zdefiniowania funkcji $\cont : K\left[ X \right] \setminus \{0\} \to K \setminus \{0\}  $.
\end{problem} 
